\documentclass[11pt,fontset=fandol]{ctexart}
\usepackage[a4paper,margin=23mm]{geometry}
\usepackage{amsmath,amssymb,mathtools,bm,mathrsfs,esint}
\usepackage{graphicx,longtable,booktabs,array,calc}
\usepackage[normalem]{ulem}
\usepackage{xcolor,hyperref}
\hypersetup{hidelinks}
\setcounter{secnumdepth}{-1}
\newcounter{none}
\providecommand{\tightlist}{\setlength{\itemsep}{0pt}\setlength{\parskip}{0pt}}
\providecommand{\passthrough}[1]{#1}
\usepackage{amssymb,mathrsfs,esint,bm,newunicodechar}
\allowdisplaybreaks
\setlength{\emergencystretch}{3em}
\providecommand{\llangle}{\langle\!\langle}
\providecommand{\rrangle}{\rangle\!\rangle}
\newunicodechar{∞}{\ensuremath{\infty}}
\newunicodechar{∈}{\ensuremath{\in}}
\newunicodechar{∑}{\ensuremath{\sum}}
\newunicodechar{∫}{\ensuremath{\int}}
\newunicodechar{≤}{\ensuremath{\leq}}
\newunicodechar{≥}{\ensuremath{\geq}}
\newunicodechar{≠}{\ensuremath{\neq}}
\newunicodechar{→}{\ensuremath{\to}}
\newunicodechar{⇒}{\ensuremath{\Rightarrow}}
\newunicodechar{α}{\ensuremath{\alpha}}
\newunicodechar{β}{\ensuremath{\beta}}
\newunicodechar{γ}{\ensuremath{\gamma}}
\newunicodechar{λ}{\ensuremath{\lambda}}
\newunicodechar{φ}{\ensuremath{\varphi}}
\newunicodechar{ε}{\ensuremath{\varepsilon}}
\newunicodechar{σ}{\ensuremath{\sigma}}
\newunicodechar{θ}{\ensuremath{\theta}}
\newunicodechar{δ}{\ensuremath{\delta}}
\newunicodechar{η}{\ensuremath{\eta}}
\newunicodechar{π}{\ensuremath{\pi}}
\newunicodechar{ξ}{\ensuremath{\xi}}
\newunicodechar{Σ}{\ensuremath{\Sigma}}
\newunicodechar{Δ}{\ensuremath{\Delta}}
\newunicodechar{∂}{\ensuremath{\partial}}
\newunicodechar{∀}{\ensuremath{\forall}}
\newunicodechar{∃}{\ensuremath{\exists}}
\newunicodechar{√}{\ensuremath{\surd}}
\newunicodechar{−}{\ensuremath{-}}
\newunicodechar{μ}{\ensuremath{\mu}}
\newunicodechar{ρ}{\ensuremath{\rho}}
\newunicodechar{τ}{\ensuremath{\tau}}
\newunicodechar{ω}{\ensuremath{\omega}}
\newunicodechar{Ω}{\ensuremath{\Omega}}
\newunicodechar{χ}{\ensuremath{\chi}}
\newunicodechar{ψ}{\ensuremath{\psi}}
\newunicodechar{Γ}{\ensuremath{\Gamma}}
\newunicodechar{Φ}{\ensuremath{\Phi}}
\newunicodechar{Λ}{\ensuremath{\Lambda}}
\newunicodechar{ℝ}{\ensuremath{\mathbb{R}}}
\newunicodechar{ℤ}{\ensuremath{\mathbb{Z}}}
\newunicodechar{ℕ}{\ensuremath{\mathbb{N}}}
\newunicodechar{△}{\ensuremath{\triangle}}
\newunicodechar{⊗}{\ensuremath{\otimes}}
\begin{document}
\section{2011年杭州电子科技大学601数学分析考研真题}\label{ux5e74ux676dux5ddeux7535ux5b50ux79d1ux6280ux5927ux5b66601ux6570ux5b66ux5206ux6790ux8003ux7814ux771fux9898}

来源：\texttt{HDU/601数学分析/2011/2011年杭州电子科技大学601数学分析考研真题.pdf}。原
PDF 页序：1, 2, 3。

\section{杭州电子科技大学}\label{ux676dux5ddeux7535ux5b50ux79d1ux6280ux5927ux5b66}

\section{2011年攻读硕士学位研究生入学考试}\label{ux5e74ux653bux8bfbux7855ux58ebux5b66ux4f4dux7814ux7a76ux751fux5165ux5b66ux8003ux8bd5}

\section{《数学分析》试题}\label{ux6570ux5b66ux5206ux6790ux8bd5ux9898}

\begin{quote}
助手整理：按同题答案及本地真题 OCR 恢复漏项和误识符号。第四大题第 3
题仍保留原文条件 \([a,b]\)；该条件与积分区间的关系见题后说明。
\end{quote}

（试题共四大题，3页，150分）

姓名 报考专业 准考证号

【所有答案必须写在答题纸上，做在试卷或草稿纸上无效！】

\subsection{一、填空题（每小题5分，共
40分）}\label{ux4e00ux586bux7a7aux9898ux6bcfux5c0fux98985ux5206ux5171-40ux5206}

\begin{enumerate}
\def\labelenumi{\arabic{enumi}.}
\tightlist
\item
  设
  \(f(x)=\begin{cases}\dfrac{\sin ax}{x},&x>0,\\2,&x=0,\\\dfrac{\ln(1-3x)}{bx},&x<0.\end{cases}\)，要使函数在
  \(x=0\) 处连续，则 \(a=\)\_\_\_\_，\(b=\)\_\_\_\_。
\end{enumerate}

2．计算积分
\(\int _ { 0 } ^ { \frac { \pi } { 2 } } \cos ^ { 2 } x d x =\)

3．设函数
\(f ( x ) = \int _ { 0 } ^ { \sqrt { \ln x } } e ^ { t ^ { 2 } } d t\)，则其导数
\(f'(x)=\)\_\_\_\_。

\begin{enumerate}
\def\labelenumi{\arabic{enumi}.}
\setcounter{enumi}{3}
\item
  设
  \(\lim_{x\to\infty}\left(\dfrac{x+a}{x+11}\right)^{x/1000}=e^2\)，则
  \(a=\)\_\_\_\_。
\item
  设 \(u(x,y,z)=\sqrt[z]{y/x}\)，求
  \(\left.\dfrac{\partial u}{\partial x}\right|_{(1,2,1)}=\)\_\_\_\_。
\end{enumerate}

6．交换二次积分
\(\int _ { 0 } ^ { 2 } d x \int _ { 0 } ^ { x } f ( x , y ) d y + \int _ { 2 } ^ { \sqrt { 8 } } d x \int _ { 0 } ^ { \sqrt { 8 - x ^ { 2 } } } f ( x , y ) d y\)
的积分顺序，则积分可表达为

\begin{enumerate}
\def\labelenumi{\arabic{enumi}.}
\setcounter{enumi}{6}
\tightlist
\item
  计算 \(\iiint_Vz\,dV=\)\_\_\_\_，其中 \(V\) 是以原点为中心、\(R\)
  为半径的上半球体。
\end{enumerate}

8.方程 \(x ^ { 5 } + x - 1 = 0\) 有 个正根.

\subsection{二、计算题（4小题，共30分）}\label{ux4e8cux8ba1ux7b97ux98984ux5c0fux9898ux517130ux5206}

1．求极限
\(\lim _ { x \to 0 } { \frac { 1 - \cos x ^ { 2 } } { x ^ { 3 } \sin x } }\)
．（本小题5分）

2．计算积分
\(\int _ { 0 } ^ { + \infty } { \frac { x } { 1 + x ^ { 4 } } } d x\)
．（本小题5分）

3．求级数
\(\sum _ { n = 0 } ^ { \infty } ( - 1 ) ^ { n } { \frac { x ^ { n + 1 } } { n + 1 } }\)
的收敛域与和函数．（本小题 10 分）

4．计算积分
\(I=\int_{C^+}(-2xe^{-x^2}\sin y)\,dx+(e^{-x^2}\cos y+x^4)\,dy\)，其中
\(C^+\) 为从点 \((1,0)\) 到点 \((-1,0)\) 的半圆周
\(y=\sqrt{1-x^2}\)，\(-1\leq x\leq1\)。（本小题10分）

\subsection{三、综合题（每小题10分，共 30
分）}\label{ux4e09ux7efcux5408ux9898ux6bcfux5c0fux989810ux5206ux5171-30-ux5206}

1．二元函数
\(f ( x , y ) = \left\{ \begin{array} { c c } { { x y \displaystyle \frac { x ^ { 2 } - y ^ { 2 } } { x ^ { 2 } + y ^ { 2 } } ~ , } } & { { x ^ { 2 } + y ^ { 2 } \ne 0 } } \\ { { 0 ~ , } } & { { x ^ { 2 } + y ^ { 2 } = 0 } } \end{array} \right. ,\)

（1）求
\({ \frac { \partial f } { \partial x } } ( x , y ) , { \frac { \partial f } { \partial y } } ( x , y )\)

(2）
\(\frac { \partial ^ { 2 } f } { \partial x \partial y } , \frac { \partial ^ { 2 } f } { \partial y \partial x }\)
在原点是否相等，求出其值.

2．判别函数项级数
\(\sum _ { n = 1 } ^ { \infty } \left( 1 - \cos { \frac { x } { n } } \right)\)
在区间 \(x \in \left[ - \delta , \delta \right] , ( \delta > 0 )\)
上的一致收敛性（需说明理由）.

\begin{enumerate}
\def\labelenumi{\arabic{enumi}.}
\setcounter{enumi}{2}
\tightlist
\item
  设
  \(f(x)=\begin{cases}x^2,&x\leq0,\\\sin x,&x>0.\end{cases}\)，（1）求函数
  \(F(x)=\int_0^xf(t)\,dt\) 的表达式；（2）讨论 \(F\) 的连续性和可导性。
\end{enumerate}

\subsection{四、证明题（每小题 10 分，共 50
分)}\label{ux56dbux8bc1ux660eux9898ux6bcfux5c0fux9898-10-ux5206ux5171-50-ux5206}

1．证明不等式
\(\sin x + \tan x > 2 x \quad \left( 0 < x < { \frac { \pi } { 2 } } \right)\)

2.设函数f(x)可微，证明：f(x)的任何两个零点之间必有
\(f ( x ) + f ^ { \prime } \left( x \right)\) 的零点.

3．已知函数f(x)在{[}a,b{]}上连续，证明：

\[
2\int_0^\alpha f(x)\left(\int_x^\alpha f(y)\,dy\right)dx=\left[\int_0^\alpha f(x)\,dx\right]^2.
\]

助手整理说明：上限 \(\alpha\) 据本地真题第 3 页恢复；原文仍写 \(f\) 在
\([a,b]\) 连续。当 \(\alpha\geq0\) 时，还需要 \(f\) 在 \([0,\alpha]\)
连续，或 \([a,b]\) 包含该积分区间，才能由所给条件合法证明上式。

\begin{enumerate}
\def\labelenumi{\arabic{enumi}.}
\setcounter{enumi}{3}
\tightlist
\item
  设 \(f(x),g(x)\) 都在区间 \(I\) 上一致连续且有界，证明
  \(F(x)=f(x)g(x)\) 也在 \(I\) 上一致连续。
\end{enumerate}

5．证明：曲面 \(x y z = a ^ { 3 } \left( a > 0 \right)\)
上任一点处的切平面和三个坐标面所围四面体的体积是一常数.
\end{document}
